% !TEX root = ../../main.tex
% C6.3 - User interface implementation (translated). The screenshots show the application's Vietnamese interface;
% screen and label names in the text are their English meaning.
\section{User Interface Implementation}
\label{sec:6.3}

\subsection{Navigation and Communication with the Server}
\label{sec:6.3.1}
\label{sec:6.3.2}

Navigation uses React Router with routes grouped by role: a user without a session goes to the login screen, and a route of another role redirects to the user's own home page. This guard only hides unusable screens; permissions are checked by the server (NFR-01).

All requests go through one API layer. It attaches the CSRF token, kept only in page memory, to every data-changing request (Section~\ref{sec:5.5.2}), and turns errors into Vietnamese messages. Server errors also show a \emph{lookup code}, the first eight characters of the request identifier, so that the matching log line can be found (Section~\ref{sec:6.4}). A 401 response clears the login state and returns to the login screen. The session is refreshed only on real user activity, shared across browser tabs, so background requests do not keep it alive.

\subsection{Rendering Person Images in the Results}
\label{sec:6.3.3}

The server renders each person image from the full frame and bounding box (Section~\ref{sec:5.3.3}). For the 3:5 result card, the crop is widened with the real surrounding scene until it fits, never cutting into the person; since it may then include other people, the surroundings are dimmed to 60\% and the person is outlined in an accent colour. The detail dialog shows the full frame with a red box. Images are sent only to users entitled to the result, marked not to be cached, and a missing frame shows an \emph{Image unavailable} icon.

\subsection{Main Screens}
\label{sec:6.3.4}

The figures show the main screens by role on the demonstration data (Section~\ref{sec:7.2}). The interface is in Vietnamese (NFR-19); screen and label names are given by their English meaning.

\paragraph{Operator.} \emph{Person search} (Figure~\ref{fig:ui-search-image}) has the query form on the left (image, cameras of the area, date range, number of results) and the ranked results on the right, each with the person image in context, camera, time and score; the shared area and date appear once above the list. In attribute search (Figure~\ref{fig:ui-search-attr}), the form shows the generated English description, so users see what is actually searched. A result opens the detail dialog (Figure~\ref{fig:ui-result}), from which it is saved into a case (Figure~\ref{fig:ui-create-case}). \emph{My cases} (Figure~\ref{fig:ui-cases}) lists the cases by status with the details of the selected one; a closed case is locked until reopened.

\begin{figure}[!htbp]
\centering
\includegraphics[width=\textwidth]{H10a-tim-kiem-anh-cat.png}
\caption{Person search by image}
\label{fig:ui-search-image}
\end{figure}

\begin{figure}[!htbp]
\centering
\includegraphics[width=\textwidth]{H10b-tim-kiem-thuoc-tinh-cat.png}
\caption{Person search by attributes}
\label{fig:ui-search-attr}
\end{figure}

\begin{figure}[!htbp]
\centering
\begin{subfigure}[b]{0.61\textwidth}
    \centering
    \includegraphics[width=\linewidth]{H10c-chi-tiet-ket-qua-cat.png}
    \caption{Details of a result}
    \label{fig:ui-result}
\end{subfigure}\hfill
\begin{subfigure}[b]{0.36\textwidth}
    \centering
    \includegraphics[width=\linewidth]{H10d-tao-vu-viec-cat.png}
    \caption{Creating a case from a result}
    \label{fig:ui-create-case}
\end{subfigure}
\caption{The result detail dialog and the case creation dialog}
\label{fig:ui-result-case}
\end{figure}

\begin{figure}[!htbp]
\centering
\includegraphics[width=\textwidth]{H10e-vu-viec-cua-toi-cat.png}
\caption{The Operator's case management screen}
\label{fig:ui-cases}
\end{figure}

\paragraph{Viewer.} \emph{Overview} (Figure~\ref{fig:ui-overview}) shows cases by status, saved results and recent cases; a case opens in \emph{Case files}, a read-only version of Figure~\ref{fig:ui-cases}.

\begin{figure}[!htbp]
\centering
\includegraphics[width=\textwidth]{H10f-tong-quan-cat.png}
\caption{The Viewer's overview screen}
\label{fig:ui-overview}
\end{figure}

\paragraph{Administrator.} \emph{Cameras} (Figure~\ref{fig:ui-cameras}) lists the cameras by area with their connection status (\emph{Uploaded video} for file-only cameras), AI processing status and actions. \emph{AI models} (Figure~\ref{fig:ui-models}) selects the Detector and Tracker, marking unusable models and showing the encoder as fixed (Section~\ref{sec:6.2.1}). \emph{Video processing} (Figure~\ref{fig:ui-video}) uploads a video with a sampling preset and follows each job by frames and ready tracks. The monitoring screens are \emph{System status} (Figure~\ref{fig:ui-status}, Section~\ref{sec:6.4}) and \emph{AI check} (Figure~\ref{fig:ui-diagnostics}, Section~\ref{sec:6.2.5}).

\begin{figure}[!htbp]
\centering
\includegraphics[width=\textwidth]{H10g-camera-cat.png}
\caption{Camera management screen}
\label{fig:ui-cameras}
\end{figure}

\begin{figure}[!htbp]
\centering
\includegraphics[width=\textwidth]{H10h-mo-hinh-ai-cat.png}
\caption{AI model configuration screen}
\label{fig:ui-models}
\end{figure}

\begin{figure}[!htbp]
\centering
\includegraphics[width=\textwidth]{H10i-xu-ly-video-cat.png}
\caption{Uploaded video processing screen}
\label{fig:ui-video}
\end{figure}

\begin{figure}[!htbp]
\centering
\includegraphics[width=\textwidth]{H10j-trang-thai-he-thong-cat.png}
\caption{System status screen}
\label{fig:ui-status}
\end{figure}

\begin{figure}[!htbp]
\centering
\includegraphics[width=\textwidth]{H10k-kiem-tra-ai-cat.png}
\caption{AI check screen}
\label{fig:ui-diagnostics}
\end{figure}
